\documentclass[../DoAn.tex]{subfiles}
\begin{document}

Chương 2 đã xác định mười nhóm chức năng cần xây dựng cùng các yêu cầu phi chức
năng về hiệu năng, bảo mật và khả năng bảo trì. Chương này trình bày các công nghệ
và nền tảng lý thuyết được lựa chọn để hiện thực hóa các yêu cầu đó, đồng thời
phân tích lý do lựa chọn so với các phương án thay thế.

\section{Kiến trúc REST API và mô hình MVC}
\label{section:3.1}

Hệ thống được xây dựng theo kiến trúc \textit{client--server} với giao tiếp qua
REST API. Phía máy chủ áp dụng mô hình MVC, tách biệt tầng
định tuyến (Route), tầng xử lý nghiệp vụ (Controller) và tầng dữ liệu (Model).
Kiến trúc này đáp ứng yêu cầu phi chức năng về khả năng bảo trì (mục
\ref{section:2.4}) và tạo nền tảng để các nhóm chức năng ở mục \ref{section:2.2}
được triển khai độc lập theo từng module.

Bảng \ref{tab:api-style} so sánh các phong cách thiết kế API phổ biến.

\begin{table}[H]
\centering
\caption{So sánh các phong cách thiết kế API}
\label{tab:api-style}
\begin{tabular}{|p{3cm}|p{3.5cm}|p{3.5cm}|p{2cm}|}
\hline
\textbf{Tiêu chí} & \textbf{REST} & \textbf{GraphQL} & \textbf{gRPC} \\
\hline
Độ phức tạp triển khai & Thấp & Trung bình & Cao \\
\hline
Hỗ trợ trình duyệt & Có sẵn & Cần thư viện & Cần proxy \\
\hline
Tích hợp \break VNPay & Tương thích tốt & Phức tạp hơn & Không phù hợp \\
\hline
Phù hợp CRUD & Tốt & Tốt & Trung bình \\
\hline
\end{tabular}
\end{table}

REST được chọn vì tương thích trực tiếp với trình duyệt, phù hợp với luồng tích
hợp thanh toán VNPay (yêu cầu redirect HTTP), và đủ đáp ứng các thao tác CRUD
của hệ thống bán lẻ mà không cần độ linh hoạt truy vấn của GraphQL.

\section{Nền tảng máy chủ: Node.js và Express.js}
\label{section:3.2}

\textbf{Node.js}~\cite{nodejs} là môi trường thực thi JavaScript phía máy chủ, sử
dụng vòng lặp sự kiện đơn luồng không chặn (non-blocking I/O), phù hợp với các
ứng dụng xử lý nhiều kết nối đồng thời. \textbf{Express.js}~\cite{expressjs} là
framework tối giản chạy trên Node.js, cung cấp hệ thống định tuyến và middleware
linh hoạt.

Hai công nghệ này đáp ứng yêu cầu xử lý đồng thời nhiều kết nối socket (chatbot
thời gian thực, mục \ref{subsection:2.2.1}) và thời gian phản hồi dưới 2 giây
(mục \ref{section:2.4}).

\begin{table}[H]
\centering
\caption{So sánh các nền tảng máy chủ}
\label{tab:backend}
\begin{tabular}{|p{3cm}|p{3cm}|p{3cm}|p{3cm}|}
\hline
\textbf{Tiêu chí} & \textbf{Node.js + \break Express} & \textbf{Django (Python)} & \textbf{Spring Boot (Java)} \\
\hline
Mô hình I/O & Bất đồng bộ & Đồng bộ (mặc định) & Đa luồng \\
\hline
Chia sẻ code với client & Có (JavaScript) & Không & Không \\
\hline
Tích hợp Socket.io & Gốc & Cần thư viện & Cần cấu hình \\
\hline
Thời gian khởi động & Nhanh & Trung bình & Chậm \\
\hline
\end{tabular}
\end{table}

Node.js được chọn vì cho phép dùng JavaScript xuyên suốt cả client và server,
giảm chi phí chuyển đổi ngữ cảnh; đồng thời tích hợp Socket.io một cách tự nhiên
để hỗ trợ chatbot thời gian thực.

\section{Cơ sở dữ liệu: MongoDB và Mongoose}
\label{section:3.3}

MongoDB~\cite{mongodb} là hệ quản trị cơ sở dữ liệu NoSQL hướng tài liệu (document-
oriented), lưu dữ liệu dưới dạng BSON. Mongoose~\cite{mongoose} là thư viện ODM
(Object Document Mapper) cung cấp schema validation và middleware cho MongoDB trên
Node.js.

MongoDB đáp ứng hai yêu cầu đặc thù xác định ở mục \ref{section:2.4}: cấu trúc
sản phẩm đa biến thể (kích thước, màu sắc) với số lượng thuộc tính không cố định,
và cấu hình phí vận chuyển.

\begin{table}[H]
\centering
\caption{So sánh các hệ quản trị cơ sở dữ liệu}
\label{tab:db}
\begin{tabular}{|p{3.5cm}|p{3cm}|p{3cm}|p{3cm}|}
\hline
\textbf{Tiêu chí} & \textbf{MongoDB} & \textbf{PostgreSQL} & \textbf{MySQL} \\
\hline
Schema linh hoạt & Có & Một phần (JSONB) & Không \\
\hline
Tài liệu lồng nhau & Tự nhiên & Cần nhiều join & Cần nhiều join \\
\hline
Tích hợp Node.js & Mongoose (ODM) & Sequelize & Sequelize \\
\hline
Phù hợp dữ liệu phân cấp & Cao & Trung bình & Thấp \\
\hline
\end{tabular}
\end{table}

MongoDB được chọn vì schema linh hoạt cho phép mỗi sản phẩm lưu danh sách biến
thể với số lượng tùy ý mà không cần thay đổi cấu trúc bảng, đồng thời cấu hình
phí vận chuyển theo địa bàn được biểu diễn tự nhiên bằng tài liệu lồng nhau.

\section{Giao diện người dùng: React.js và các thư viện}
\label{section:3.4}

React~\cite{react} là thư viện JavaScript do Meta phát triển, xây dựng giao diện
theo mô hình component tái sử dụng với cơ chế Virtual DOM để cập nhật giao diện
hiệu quả. Đồ án sử dụng thêm:

\begin{itemize}
  \item \textbf{React Router DOM v6}~\cite{reactrouter}: định tuyến phía client,
        hỗ trợ SPA \break (Single Page Application) cho cả giao diện mua sắm và trang
        quản trị.
  \item \textbf{Redux Toolkit}~\cite{redux}: quản lý trạng thái toàn cục (giỏ
        hàng, thông tin người dùng đăng nhập) nhất quán giữa các component.
  \item \textbf{Bootstrap 5 + React-Bootstrap}: đảm bảo giao diện đáp ứng
        (responsive) theo yêu cầu tương thích đa thiết bị (mục \ref{section:2.4}).
  \item \textbf{Chart.js / Recharts}: trực quan hóa dữ liệu doanh thu trên trang
        quản trị.
\end{itemize}

\begin{table}[H]
\centering
\caption{So sánh các framework/thư viện giao diện}
\label{tab:frontend}
\begin{tabular}{|p{3.5cm}|p{3cm}|p{3cm}|p{3cm}|}
\hline
\textbf{Tiêu chí} & \textbf{React} & \textbf{Vue.js} & \textbf{Angular} \\
\hline
Mô hình & Thư viện UI & Framework & Framework đầy đủ \\
\hline
Độ linh hoạt & Cao & Cao & Thấp hơn \\
\hline
Hệ sinh thái & Rất lớn & Lớn & Lớn \\
\hline
Đường cong học tập & Trung bình & Thấp & Cao \\
\hline
Quản lý trạng thái & Redux Toolkit & Pinia/Vuex & NgRx/Service \\
\hline
\end{tabular}
\end{table}

React được chọn vì hệ sinh thái phong phú (Redux Toolkit, React Router, thư viện
biểu đồ) cho phép xây dựng cả storefront lẫn trang quản trị trong cùng một dự
án, đồng thời mô hình component tái sử dụng phù hợp với giao diện có nhiều phần
tử lặp lại như danh sách sản phẩm, thẻ đơn hàng.

\section{Xác thực và phân quyền}
\label{section:3.5}

\subsection{JSON Web Token}
\label{subsection:3.5.1}

JSON Web Token (JWT)~\cite{rfc7519} là tiêu chuẩn mở (RFC 7519) định nghĩa cách
truyền thông tin an toàn giữa các bên dưới dạng đối tượng JSON được ký số. Token
gồm ba phần: Header, Payload và Signature, được ký bằng thuật toán HMAC-SHA256.

JWT đáp ứng yêu cầu xác thực người dùng ở mục \ref{section:2.4}: server không cần
lưu trạng thái phiên (stateless), mỗi yêu cầu API kèm token để kiểm tra quyền
trước khi xử lý.

\begin{table}[H]
\centering
\caption{So sánh các cơ chế xác thực}
\label{tab:auth}
\begin{tabular}{|p{3.5cm}|p{3.2cm}|p{3.2cm}|p{2.8cm}|}
\hline
\textbf{Tiêu chí} & \textbf{JWT} & \textbf{Session--Cookie} & \textbf{OAuth2} \\
\hline
Lưu trạng thái server & Không (stateless) & Có (session store) & Có thể có hoặc không \\
\hline
Phù hợp SPA / REST & Tốt & Cần cấu hình CORS & Phức tạp hơn \\
\hline
Kiểm soát hết hạn & Có (exp claim) & Có & Có \\
\hline
Triển khai & Đơn giản & Trung bình & Phức tạp \\
\hline
\end{tabular}
\end{table}

\subsection{Mã hóa mật khẩu với bcrypt}
\label{subsection:3.5.2}

bcrypt~\cite{provos1999} là thuật toán băm mật khẩu được thiết kế có chủ ý chậm,
kết hợp salt ngẫu nhiên để chống tấn công từ điển và rainbow table. Mật khẩu
người dùng được băm trước khi lưu vào MongoDB, đáp ứng yêu cầu bảo mật ở mục
\ref{section:2.4}. So với SHA-256, bcrypt an toàn hơn vì chi phí tính toán có thể
điều chỉnh theo khả năng phần cứng; so với Argon2, bcrypt có thư viện hỗ trợ
Node.js tốt hơn \break (\texttt{bcrypt} v5).

\subsection{Phân quyền dựa trên vai trò (RBAC)}
\label{subsection:3.5.3}

Mô hình RBAC (Role-Based Access Control)~\cite{sandhu1996} gán quyền cho vai trò
thay vì gán trực tiếp cho người dùng; người dùng được gán một hoặc nhiều vai trò.
Mô hình này hiện thực hóa yêu cầu phân quyền nhân viên theo vai trò với lịch sử tác động hệ thống (mục \ref{subsection:2.3.4}).

Hệ thống định nghĩa ba vai trò chính (\textit{admin}, \textit{user},
\textit{staff}) với bộ quyền hành động (Action) tương ứng lưu trong MongoDB.
Middleware kiểm tra vai trò tại tầng API trước mỗi yêu cầu. So với ACL (Access
Control List) gán quyền theo từng người dùng, RBAC dễ bảo trì hơn khi số lượng
nhân viên tăng; so với ABAC (Attribute-Based), RBAC đơn giản hơn và đủ đáp ứng
yêu cầu của hệ thống bán lẻ quy mô vừa và nhỏ.

\section{Giao tiếp thời gian thực: Socket.io}
\label{section:3.6}

Socket.io~\cite{socketio} là thư viện cho phép giao tiếp hai chiều theo thời gian
thực giữa client và server thông qua \break WebSocket (với fallback về long-polling khi
WebSocket không khả dụng). Thư viện này hiện thực hóa chức năng hỗ trợ khách hàng
qua chatbot thời gian thực (mục \ref{subsection:2.2.1}).

\begin{table}[H]
\centering
\caption{So sánh các công nghệ giao tiếp thời gian thực}
\label{tab:realtime}
\begin{tabular}{|p{3.5cm}|p{3cm}|p{3cm}|p{3cm}|}
\hline
\textbf{Tiêu chí} & \textbf{Socket.io} & \textbf{WebSocket thuần} & \textbf{SSE} \\
\hline
Hướng giao tiếp & Hai chiều & Hai chiều & Một chiều \\
\hline
Tự động reconnect & Có & Không (tự làm) & Có \\
\hline
Fallback & Có \break (long-polling) & Không & Không \\
\hline
Tích hợp Node.js & Tích hợp sẵn & Thư viện \texttt{ws} & HTTP gốc \\
\hline
\end{tabular}
\end{table}

Socket.io được chọn vì cơ chế tự động kết nối lại và fallback giúp đảm bảo độ
tin cậy kết nối theo yêu cầu ở mục \ref{section:2.4}, đồng thời socket.io đã được  tích hợp 
với Node.js không cần cấu hình thêm.

\section{Tích hợp thanh toán: VNPay}
\label{section:3.7}

VNPay~\cite{vnpay} là cổng thanh toán trực tuyến phổ biến tại Việt Nam, hỗ trợ
thanh toán qua thẻ nội địa ATM, thẻ quốc tế và ví điện tử. Tích hợp thực hiện
theo luồng redirect: khi khách hàng chọn thanh toán VNPay, server tạo URL thanh
toán có chữ ký HMAC-SHA512, chuyển hướng trình duyệt đến cổng VNPay; sau khi
hoàn tất, VNPay gọi callback URL để xác nhận kết quả. Luồng này khớp với quy
trình đặt hàng và thanh toán mô tả ở mục \ref{subsection:2.2.7}.

So với các lựa chọn thay thế: MoMo và ZaloPay cũng là cổng thanh toán phổ biến
tại Việt Nam nhưng yêu cầu xác minh doanh nghiệp phức tạp hơn để sử dụng môi
trường sandbox; Stripe không hỗ trợ thẻ ATM nội địa Việt Nam. VNPay cung cấp môi
trường sandbox đầy đủ cho mục đích phát triển và phù hợp với đối tượng người dùng
trong nước của hệ thống ALIBARBIE.

\section{Thông báo qua email: Nodemailer và Google SMTP}
\label{section:3.8}

Nodemailer~\cite{nodemailer} là module Node.js để gửi email, được cấu hình sử dụng
Google SMTP thông qua OAuth2 (googleapis). Chức năng này đáp ứng yêu cầu gửi email
xác nhận đơn hàng và thông báo kết quả hoàn trả (mục \ref{subsection:2.3.1},
\ref{subsection:2.3.2}).

So với SendGrid hay Amazon SES, Google SMTP phù hợp hơn cho giai đoạn phát triển
vì không yêu cầu đăng ký tài khoản dịch vụ trả phí; so với gửi SMTP trực tiếp
bằng tên và mật khẩu, OAuth2 an toàn hơn vì không lưu mật khẩu trong mã nguồn.

Bảng \ref{tab:tech-summary} tóm tắt sự tương ứng giữa các công nghệ và yêu cầu
chức năng đã xác định ở Chương 2.

\begin{table}[H]
\centering
\caption{Tóm tắt công nghệ và yêu cầu chức năng tương ứng}
\label{tab:tech-summary}
\begin{tabular}{|p{4cm}|p{8.5cm}|}
\hline
\textbf{Công nghệ} & \textbf{Yêu cầu chức năng được giải quyết} \\
\hline
Node.js + Express.js & Nền tảng API phục vụ toàn bộ chức năng; xử lý đồng thời socket \\
\hline
MongoDB  \break + Mongoose & Sản phẩm đa biến thể; cấu hình phí vận chuyển phân cấp \\
\hline
React + Redux Toolkit & Giao diện mua sắm và quản trị; quản lý trạng thái giỏ hàng \\
\hline
JWT + bcrypt & Xác thực người dùng; bảo mật mật khẩu \\
\hline
RBAC & Phân quyền nhân viên; nhật ký kiểm toán hành động \\
\hline
Socket.io & Chatbot hỗ trợ khách hàng thời gian thực \\
\hline
VNPay & Thanh toán trực tuyến trong quy trình đặt hàng \\
\hline
Nodemailer + Google SMTP & Email xác nhận đơn hàng; thông báo kết quả hoàn trả \\
\hline
\end{tabular}
\end{table}

Chương này đã trình bày tám nhóm công nghệ được lựa chọn cùng lý do so với các
phương án thay thế, gắn với các yêu cầu cụ thể đã phân tích ở Chương 2. Chương
tiếp theo sẽ trình bày thiết kế và hiện thực hóa hệ thống dựa trên nền tảng công
nghệ này.

%%%%%%%%%%%%%%%%%%%%%%%%%%%%%%%%%%%
\end{document}